\chapter{Auditoria do Setor Público}

\section{Auditoria pública e accountability}

Auditoria do setor público é processo sistemático de obtenção e avaliação objetiva de evidências para determinar se determinado objeto está em conformidade com critérios aplicáveis ou apresenta desempenho esperado, comunicando resultados aos usuários previstos.

Sua função institucional é ampliar transparência, accountability e confiança na administração de recursos públicos.

A lógica de um trabalho pode ser resumida em:
\[
\text{mandato}
\rightarrow
\text{objetivo}
\rightarrow
\text{questões}
\rightarrow
\text{critérios}
\rightarrow
\text{riscos}
\rightarrow
\text{procedimentos}
\rightarrow
\text{evidências}
\rightarrow
\text{achados}
\rightarrow
\text{conclusão}
\rightarrow
\text{monitoramento}.
\]

Conclusão sem cadeia verificável é opinião, não resultado de auditoria adequadamente fundamentado.

\section{Princípios fundamentais — ISSAI e NBASP}

As normas internacionais das Entidades Fiscalizadoras Superiores (INTOSAI) e as Normas Brasileiras de Auditoria do Setor Público estruturam princípios fundamentais aplicáveis ao trabalho de auditoria pública.

Elementos conceituais incluem:
\begin{itemize}
\item auditor;
\item parte responsável;
\item usuários previstos;
\item objeto;
\item critérios;
\item evidência;
\item relatório/conclusão.
\end{itemize}

\subsection{Objeto e critério}

\termo{Objeto} é aquilo que será avaliado: demonstrações, contratação, política, processo, sistema, controle, programa ou conjunto de transações.

\termo{Critério} é referência usada para comparar o objeto. Pode decorrer de lei, contrato, norma, política, padrão técnico ou meta legitimamente definida.

Sem critério adequado não existe base objetiva para classificar condição como regular, irregular, eficiente ou deficiente.

\section{Trabalhos de certificação e relatório direto}

Em trabalho de \termo{certificação}, parte responsável mede ou apresenta informação do objeto e o auditor obtém evidência sobre essa informação.

Em \termo{relatório direto}, o auditor mede ou avalia diretamente o objeto segundo critérios e apresenta o resultado.

A distinção é importante porque altera a relação entre objeto, informação do objeto e conclusão.

\section{Tipos de auditoria}

\subsection{Auditoria financeira}

Busca determinar se informação financeira é apresentada, em todos os aspectos relevantes, conforme estrutura de relatório aplicável.

Trabalha intensamente com materialidade, risco de distorção e evidência sobre saldos, transações e divulgações.

\subsection{Auditoria de conformidade}

Avalia se atividades, transações, informações ou decisões obedecem autoridades aplicáveis, como legislação, regulamentos, contratos e políticas.

\subsection{Auditoria operacional}

Examina economicidade, eficiência, eficácia e efetividade de programas, organizações ou processos.

\begin{table}[ht]
\centering
\small
\begin{tabularx}{\textwidth}{l X}
\toprule
Dimensão & Pergunta \\
\midrule
Economicidade & recursos foram obtidos em custo e condições adequados? \\
Eficiência & recursos foram convertidos em produtos com boa produtividade? \\
Eficácia & metas/produtos previstos foram alcançados? \\
Efetividade & resultados produziram impacto real pretendido? \\
\bottomrule
\end{tabularx}
\end{table}

Uma política pode ser eficaz ao entregar 100\% das vagas previstas e pouco efetiva se não melhorar o indicador social que justificou sua existência.

\section{Independência e objetividade}

Independência reduz interferência que comprometa julgamento. Objetividade exige avaliação imparcial da evidência.

Ameaças incluem:
\begin{itemize}
\item interesse próprio;
\item autorrevisão;
\item familiaridade;
\item intimidação;
\item defesa de posição da unidade auditada;
\item participação prévia em decisão gerencial do objeto.
\end{itemize}

Salvaguardas podem incluir substituição de membro, revisão independente, transparência e segregação de funções.

\section{Ética e ceticismo profissional}

Ceticismo profissional é postura questionadora e atenção a condições que possam indicar erro, fraude ou evidência inconsistente.

Não significa presumir má-fé. O auditor deve estar disposto a aceitar explicação sustentada por evidência e a rejeitar explicação não corroborada.

\section{Julgamento profissional}

Normas não eliminam necessidade de julgamento. O auditor decide sobre:
\begin{itemize}
\item materialidade;
\item risco;
\item extensão dos testes;
\item amostra;
\item confiabilidade da evidência;
\item relevância de achado;
\item forma da conclusão.
\end{itemize}

Julgamento profissional precisa ser documentado de forma que outro auditor experiente compreenda premissas e raciocínio.

\section{Planejamento de auditoria}

Planejamento transforma mandato amplo em trabalho executável.

\subsection{Conhecimento do objeto}

O auditor deve compreender:
\begin{itemize}
\item objetivos institucionais;
\item processos;
\item legislação;
\item sistemas;
\item dados;
\item controles;
\item fornecedores;
\item responsáveis;
\item histórico de problemas;
\item ambiente externo.
\end{itemize}

\subsection{Objetivo e escopo}

Objetivo descreve o que será concluído. Escopo delimita período, unidades, localidades, sistemas e temas.

Escopo amplo demais produz trabalho superficial; estreito demais pode omitir riscos essenciais.

\section{Questões de auditoria}

Questões devem ser claras e respondíveis.

Fraca: ``A segurança do sistema é adequada?''

Melhor: ``Os acessos privilegiados são concedidos, revisados e revogados conforme critérios estabelecidos e de modo a preservar segregação e rastreabilidade?''

Subquestões podem decompor problemas complexos.

\section{Critérios de auditoria}

Critérios devem ser relevantes, completos o suficiente, confiáveis, neutros e compreensíveis segundo o tipo de trabalho.

Fontes:
\begin{itemize}
\item Constituição e leis;
\item regulamentos;
\item contratos;
\item políticas aprovadas;
\item normas técnicas;
\item metas;
\item benchmarks tecnicamente justificáveis.
\end{itemize}

Benchmark de mercado não se converte automaticamente em obrigação. É necessário demonstrar aplicabilidade ao objeto.

\section{Risco de auditoria}

Na auditoria financeira, modelo clássico:
\[
RA=RI\times RC\times RD,
\]
onde:
\begin{itemize}
\item $RI$: risco inerente;
\item $RC$: risco de controle;
\item $RD$: risco de detecção;
\item $RA$: risco de auditoria.
\end{itemize}

Risco inerente e de controle pertencem ao objeto/entidade; risco de detecção é influenciado por procedimentos do auditor.

Se RI e RC são altos e RA desejado é baixo, o auditor precisa reduzir RD com procedimentos mais fortes ou extensos.

\begin{exemplo}[efeito no risco de detecção]
Suponha meta didática $RA=5\%$, $RI=80\%$ e $RC=50\%$:
\[
RD=\frac{0{,}05}{0{,}8\times0{,}5}=0{,}125=12{,}5\%.
\]
Isso ilustra que riscos altos no objeto exigem menor tolerância a falha dos testes.
\end{exemplo}

O modelo é simplificação quantitativa; na prática, avaliações são frequentemente qualitativas e baseadas em normas específicas.

\section{Materialidade}

Materialidade é magnitude ou natureza de erro/irregularidade que pode influenciar decisão dos usuários previstos.

Pode ser:
\begin{itemize}
\item quantitativa;
\item qualitativa;
\item contextual.
\end{itemize}

Vazamento de base de denunciantes pode ser altamente material mesmo com baixo impacto financeiro.

\section{Materialidade de desempenho}

Em auditoria financeira, pode-se definir materialidade de desempenho abaixo da materialidade global para reduzir probabilidade de conjunto de distorções não detectadas exceder materialidade.

O raciocínio é reservar margem para erros distribuídos em múltiplas áreas.

\section{Risco de fraude}

Fraude envolve ato intencional para obter vantagem indevida ou produzir representação enganosa.

Triângulo da fraude clássico:
\begin{itemize}
\item pressão/incentivo;
\item oportunidade;
\item racionalização.
\end{itemize}

Modelos ampliados incluem capacidade.

O auditor não é garantidor de inexistência de fraude, mas deve considerar riscos e responder conforme seu mandato e normas.

\section{Controles internos}

Controle interno é processo desenhado para fornecer segurança razoável quanto a objetivos.

O COSO Internal Control trabalha com cinco componentes:
\begin{enumerate}
\item ambiente de controle;
\item avaliação de riscos;
\item atividades de controle;
\item informação e comunicação;
\item monitoramento.
\end{enumerate}

Controles podem ser preventivos, detectivos ou corretivos e manuais ou automatizados.

\section{Desenho, implementação e efetividade operacional}

Um controle pode falhar em três níveis:
\begin{itemize}
\item \textbf{desenho}: mesmo executado, não reduz risco adequadamente;
\item \textbf{implementação}: foi definido, mas não colocado em funcionamento;
\item \textbf{efetividade operacional}: existe e foi implantado, mas não opera consistentemente.
\end{itemize}

\begin{exemplo}[recertificação de acessos]
Desenho: política exige revisão trimestral por gestor adequado.\\
Implementação: ferramenta envia campanhas trimestrais.\\
Operação: gestores revisam efetivamente, removem acessos e há evidência íntegra.
\end{exemplo}

\section{Matriz de planejamento}

\begin{table}[ht]
\centering
\small
\begin{tabularx}{\textwidth}{X X X X}
\toprule
Questão & Critério & Evidência necessária & Procedimento \\
\midrule
SLA foi cumprido? & contrato/TR & logs e tickets & recalcular disponibilidade e reconciliar com pagamento. \\
Privilégios são revisados? & política/IAM & população, aprovações, logs & testar desenho e amostra/censo de recertificações. \\
Backup atende RPO/RTO? & BCP/contrato & jobs, PITR, testes & inspecionar e observar/reexecutar restore. \\
\bottomrule
\end{tabularx}
\end{table}

Procedimento sem relação com questão produz evidência inútil.

\section{Procedimentos de auditoria}

\subsection{Inspeção}

Exame de registros, documentos e ativos.

\subsection{Observação}

Acompanhar execução de procedimento. Limitação: pessoas podem mudar comportamento por saber que estão sendo observadas.

\subsection{Indagação}

Obter explicações. Normalmente precisa de corroboração para questões relevantes.

\subsection{Confirmação externa}

Obter resposta diretamente de terceira parte independente.

\subsection{Recálculo}

Verificar exatidão matemática ou lógica.

\subsection{Reexecução}

Executar novamente procedimento/controle de forma independente.

\subsection{Procedimentos analíticos}

Comparar relações, tendências, razões e expectativas para identificar inconsistências.

\section{Teste de controle versus procedimento substantivo}

\termo{Teste de controle} verifica desenho e/ou operação de controle.

\termo{Procedimento substantivo} busca detectar diretamente erro, distorção ou irregularidade na informação/objeto.

\begin{exemplo}[pagamento]
Controle: aprovação eletrônica antes da liquidação. Teste de controle: verificar se aprovações ocorreram de modo adequado em amostra.

Substantivo: recalcular pagamentos e confrontar com contrato/entregas para identificar sobrepreço ou quantidade indevida.
\end{exemplo}

\section{Evidência de auditoria}

Evidência deve ser \termo{suficiente} em quantidade e \termo{apropriada} em qualidade.

Apropriação envolve relevância e confiabilidade.

\subsection{Confiabilidade}

Tende a ser maior quando:
\begin{itemize}
\item fonte é externa e independente;
\item obtida diretamente pelo auditor;
\item produzida sob controles confiáveis;
\item documental/eletrônica íntegra;
\item original ou verificável.
\end{itemize}

Essas são tendências, não hierarquia absoluta. Confirmação externa pode ser falsa; log interno pode ser altamente confiável se protegido e validado.

\section{Contradição entre evidências}

Se duas fontes divergem, auditor não deve escolher a mais conveniente. Deve investigar causa.

Exemplo: ferramenta de monitoramento registra 99,99\%, tickets indicam queda de 4 h. Possibilidades:
\begin{itemize}
\item monitor mede endpoint inadequado;
\item tickets incluem indisponibilidade parcial;
\item relógios divergentes;
\item exclusões do SLA;
\item manipulação de dados.
\end{itemize}

A inconsistência é evidência que exige procedimento adicional.

\section{Evidência digital}

Dados digitais precisam de cadeia de origem, integridade, completude e transformação.

\subsection{Hash}

Hash ajuda a demonstrar que cópia não mudou após coleta.

Não prova que fonte original não foi alterada antes.

\subsection{Metadados}

Timestamp, usuário, sistema, timezone e versão são essenciais para interpretar evento.

\subsection{Cadeia de custódia}

Documenta quem coletou, quando, como armazenou, quem acessou e como transferiu evidência.

\section{Reprodutibilidade de analytics}

Se relatório afirma ``2.341 pagamentos duplicados'', deve ser possível reproduzir:
\begin{itemize}
\item query/script;
\item snapshot da base;
\item filtros;
\item tratamento de NULL;
\item critérios de duplicidade;
\item versão de dependências;
\item resultado intermediário.
\end{itemize}

Dashboard sem query e lineage é evidência frágil.

\section{Amostragem em auditoria}

Amostragem examina menos de 100\% da população para obter conclusão sobre conjunto.

\subsection{Amostragem estatística}

Usa seleção probabilística e teoria estatística para avaliar risco de amostragem.

\subsection{Amostragem não estatística}

Usa julgamento profissional sem inferência estatística formal. Ainda deve ser planejada e documentada.

\section{Risco de amostragem}

É risco de conclusão baseada na amostra diferir da conclusão que seria obtida examinando toda população.

Aumentar amostra reduz risco, mas custo cresce.

Tamanho depende de:
\begin{itemize}
\item nível de confiança;
\item desvio tolerável;
\item desvio esperado;
\item materialidade;
\item variabilidade;
\item desenho de seleção.
\end{itemize}

\section{Amostragem aleatória e sistemática}

Aleatória simples atribui chance apropriada a cada item.

Sistemática seleciona a cada intervalo $k=N/n$, após início aleatório.

Se população possui periodicidade relacionada a $k$, seleção pode ser enviesada.

\section{Amostragem estratificada}

Divide população em estratos mais homogêneos ou de risco diferente.

Exemplo: examinar 100\% dos pagamentos acima de R\$ 5 milhões e amostrar estatisticamente os demais. Isso combina cobertura de itens materialmente relevantes com eficiência.

\section{Monetary Unit Sampling}

Em amostragem por unidade monetária, unidades de moeda possuem chance de seleção, fazendo itens de maior valor terem maior probabilidade.

É útil em testes de superavaliação, mas possui limitações para saldos negativos e subavaliações.

\section{Censo e analytics}

Processar 100\% da base não elimina risco de auditoria.

Ainda podem existir:
\begin{itemize}
\item dados incompletos;
\item regra analítica incorreta;
\item sistema-fonte não confiável;
\item registros fora da base;
\item falsos positivos/negativos.
\end{itemize}

Censo elimina risco de amostragem sobre aquela população disponível, não outros riscos.

\section{Achados de auditoria}

Estrutura:
\begin{itemize}
\item critério;
\item condição;
\item causa;
\item efeito/risco;
\item evidência;
\item encaminhamento.
\end{itemize}

\subsection{Condição}

Deve ser factual e quantificada quando possível.

Fraco: ``há muitas contas antigas''.

Forte: ``37 de 214 contas privilegiadas estavam sem uso há mais de 180 dias e 12 pertenciam a usuários desligados''.

\subsection{Causa}

Deve explicar mecanismo que produziu condição.

``Falha humana'' raramente é causa-raiz suficiente. Pergunte por que erro não foi prevenido/detectado.

\subsection{Efeito}

Pode ser dano ocorrido ou risco plausível. Não se deve afirmar dano financeiro sem quantificação/evidência.

\section{Matriz de achados}

Uma matriz organiza:
\begin{itemize}
\item situação encontrada;
\item critério;
\item evidência;
\item causa;
\item efeito;
\item boas práticas;
\item proposta de encaminhamento;
\item benefício esperado.
\end{itemize}

Ela é ferramenta de raciocínio; não substitui texto do relatório.

\section{Recomendações e determinações}

Recomendação deve atacar causa e preservar espaço gerencial legítimo.

Determinação exige obrigação jurídica e fundamento para impor providência.

\begin{example}[recomendação orientada a resultado]
Em vez de ``compre ferramenta X'', formular: ``implementar processo automatizado de revogação de acessos integrado ao cadastro de pessoal, com confirmação de desligamentos em prazo definido e trilha auditável''.
\end{example}

A solução específica pode ser escolhida pelo gestor, salvo requisito técnico/jurídico que a determine.

\section{Relatório de auditoria}

Qualidades:
\begin{itemize}
\item clareza;
\item objetividade;
\item completude proporcional;
\item precisão;
\item tempestividade;
\item equilíbrio;
\item rastreabilidade.
\end{itemize}

Conclusões devem refletir escopo. Auditar uma unidade não permite afirmar que ``todo o Tribunal'' possui o mesmo problema sem base adicional.

\section{Contraditório}

Manifestação da unidade pode:
\begin{itemize}
\item explicar contexto;
\item apresentar evidência nova;
\item contestar critério;
\item corrigir erro factual;
\item demonstrar ação já realizada.
\end{itemize}

Auditor deve reavaliar de modo objetivo. Manter achado apenas porque já foi escrito viola ceticismo profissional.

\section{Monitoramento}

Monitoramento examina implementação e, quando pertinente, resultado.

Estados possíveis podem distinguir:
\begin{itemize}
\item implementada;
\item parcialmente implementada;
\item não implementada;
\item não mais aplicável;
\item substituída por solução equivalente.
\end{itemize}

Evidência de ação não é necessariamente evidência de efetividade.

\section{Controle Geral de TI — ITGC}

General IT Controls sustentam ambiente tecnológico e confiabilidade de controles automatizados.

Áreas clássicas:
\begin{itemize}
\item gestão de acessos;
\item gestão de mudanças;
\item operações;
\item desenvolvimento/aquisição;
\item backup/continuidade;
\item segurança física/lógica.
\end{itemize}

Se ITGC é fraco, confiança em controles automatizados dependentes pode ser reduzida.

\section{Auditoria de acessos}

Objetivos:
\begin{itemize}
\item usuários autorizados;
\item privilégio mínimo;
\item segregação;
\item contas privilegiadas;
\item offboarding;
\item recertificação;
\item MFA;
\item logging.
\end{itemize}

\subsection{Teste de população}

Extrair todas as contas, reconciliar com RH/identidade e identificar:
\begin{itemize}
\item desligados ativos;
\item contas sem owner;
\item contas compartilhadas;
\item privilégios incompatíveis;
\item inativas;
\item contas de serviço interativas.
\end{itemize}

A regra analítica deve tratar exceções legítimas.

\section{Auditoria de mudanças}

Cadeia esperada:
\[
\text{solicitação}\rightarrow
\text{avaliação de risco}\rightarrow
\text{aprovação}\rightarrow
\text{teste}\rightarrow
\text{implantação}\rightarrow
\text{validação}\rightarrow
\text{rollback/revisão}.
\]

Testes:
\begin{itemize}
\item mudanças em produção possuem ticket?;
\item aprovador é independente?;
\item commit/artifact correspondem à mudança aprovada?;
\item emergenciais têm revisão posterior?;
\item desenvolvedor pode alterar produção diretamente?
\end{itemize}

\section{Auditoria de operações}

Abrange jobs, batch, monitoring, incidentes, capacidade, backup e disponibilidade.

Uma falha de job sem alerta pode ser mais relevante que um job que falha visivelmente, porque produz dados desatualizados silenciosamente.

\section{Controles de aplicação}

Controles de aplicação atuam em transações e regras do sistema.

Categorias:
\begin{itemize}
\item entrada;
\item processamento;
\item saída;
\item interface;
\item autorização;
\item trilha de auditoria.
\end{itemize}

\section{Controles de entrada}

Exemplos:
\begin{itemize}
\item formato;
\item faixa;
\item obrigatoriedade;
\item duplicidade;
\item check digit;
\item autorização.
\end{itemize}

Validar apenas front-end não é suficiente; API/backend deve impor regra.

\section{Controles de processamento}

Exemplos:
\begin{itemize}
\item totais de controle;
\item reconciliação;
\item regras de cálculo;
\item transação atômica;
\item detecção de duplicatas;
\item sequenciamento.
\end{itemize}

\section{Controles de interface}

Integrações devem assegurar completude e unicidade.

Mecanismos:
\begin{itemize}
\item contagem de registros;
\item total monetário;
\item checksum;
\item identificador idempotente;
\item dead-letter queue;
\item reconciliação origem-destino.
\end{itemize}

API retornar HTTP 200 não prova que todos os registros foram processados corretamente.

\section{Auditoria de banco de dados}

Verificar:
\begin{itemize}
\item privilégios;
\item contas administrativas;
\item criptografia;
\item backup/PITR;
\item mudanças de schema;
\item logs;
\item replicação;
\item integridade;
\item segregação entre aplicação e DBA.
\end{itemize}

\section{Auditoria de segurança}

Pode usar critérios de NIST, ISO, CIS ou políticas internas conforme aplicabilidade.

Procedimentos:
\begin{itemize}
\item configuração;
\item vulnerabilidades;
\item patches;
\item IAM;
\item segmentação;
\item SIEM;
\item incident response;
\item backup;
\item testes de recuperação.
\end{itemize}

Scanner de vulnerabilidade é evidência técnica, mas falsos positivos/negativos precisam de validação.

\section{Auditoria de nuvem}

Objetos:
\begin{itemize}
\item contas/projetos;
\item IAM;
\item logging;
\item redes;
\item storage público;
\item KMS;
\item snapshots;
\item regiões;
\item custos;
\item backup;
\item CSPM/policy as code.
\end{itemize}

Evidência pode ser extraída por APIs do provedor e preservada como snapshot estruturado, preferível a dezenas de screenshots.

\section{Auditoria de contratos de TI}

A auditoria conecta:
\[
\text{necessidade}\rightarrow
\text{quantidade}\rightarrow
\text{preço}\rightarrow
\text{entrega}\rightarrow
\text{medição}\rightarrow
\text{pagamento}.
\]

Procedimentos:
\begin{itemize}
\item recalcular SLA;
\item reconciliar faturas com inventário;
\item comparar consumo licenciado;
\item analisar preços;
\item verificar aceite;
\item testar multas/glosas;
\item examinar lock-in e renovação.
\end{itemize}

\section{CAATs e auditoria assistida por computador}

Computer-Assisted Audit Techniques incluem consultas, scripts, análise de logs, amostragem automatizada e teste de controles.

Benefícios:
\begin{itemize}
\item examinar populações maiores;
\item reproduzir cálculos;
\item cruzar fontes;
\item identificar padrões raros;
\item documentar regras.
\end{itemize}

Riscos:
\begin{itemize}
\item query errada em escala;
\item dados incompletos;
\item duplicação por join;
\item interpretação inadequada;
\item alteração não versionada do código.
\end{itemize}

\section{Auditoria contínua}

Auditoria contínua usa testes automatizados/recorrentes para detectar exceções com menor latência.

Não significa que toda exceção seja achado definitivo. Regras geram alertas que precisam de triagem e contexto.

\section{Indicadores e thresholds de auditoria contínua}

Exemplo:
\begin{itemize}
\item contas de desligados ativas $>24$ h;
\item pagamentos duplicados pelo mesmo documento;
\item SLA abaixo de 99,95\%;
\item mudança emergencial sem revisão em 5 dias;
\item bucket público não autorizado.
\end{itemize}

Threshold deve equilibrar cobertura e falso positivo.

\section{Auditoria de modelos e IA}

Auditar IA requer examinar:
\begin{itemize}
\item finalidade;
\item dados;
\item leakage;
\item métricas;
\item grupos;
\item versão;
\item drift;
\item supervisão;
\item explicabilidade;
\item segurança;
\item fornecedores;
\item logs de decisão.
\end{itemize}

Boa acurácia global não é suficiente se subgrupos apresentam erro inaceitável ou se dados usados não estavam disponíveis no momento real.

\section{Fraude e análise de dados}

Red flags não são prova.

Exemplos:
\begin{itemize}
\item pagamentos próximos ao limite de autorização;
\item fornecedores compartilhando endereço/telefone;
\item sequência de notas em curto período;
\item concentração incomum;
\item valores arredondados repetitivos;
\item alteração de cadastro antes de pagamento.
\end{itemize}

Cada sinal precisa de investigação e evidência corroborativa.

\section{Qualidade da auditoria}

Qualidade envolve políticas e revisões que assegurem conformidade com normas e relatórios apropriados.

Mecanismos:
\begin{itemize}
\item supervisão;
\item revisão de papéis;
\item consultas técnicas;
\item independência;
\item competência;
\item revisão de qualidade para trabalhos relevantes;
\item lições aprendidas.
\end{itemize}

\section{Síntese conceitual}

\mapafuncional{Auditoria do Setor Público}
{Auditoria confiável = questão clara + critério adequado + procedimento proporcional + evidência suficiente e apropriada}
{Fundamentos $\rightarrow$ mandato, partes, objeto e critérios.\\Planejamento $\rightarrow$ risco, materialidade, questões e matriz.\\Execução $\rightarrow$ testes, amostragem e evidência.\\Achado $\rightarrow$ critério, condição, causa e efeito.\\TI $\rightarrow$ ITGC, application controls, dados, cloud e IA.\\Fechamento $\rightarrow$ relatório, contraditório e monitoramento.}
{Risco alto exige procedimentos mais fortes.\\Censo elimina risco de amostragem, não risco de dados.\\Controle documentado precisa operar.\\Hash prova integridade da cópia, não verdade da fonte.\\Analytics gera indício; responsabilização exige corroboração.}
{Indagação $\neq$ evidência suficiente.\\Screenshot $\neq$ população completa.\\Materialidade $\neq$ apenas valor.\\Erro humano $\neq$ causa-raiz automática.\\100\% dos dados $\neq$ 100\% de segurança.\\Recomendação $\neq$ assumir gestão.}
{Qual conclusão precisa ser sustentada?\\Qual critério se aplica?\\A população está completa?\\O procedimento responde à questão?\\A evidência é independente/reprodutível?\\A causa foi demonstrada?\\A ação corrigiu o efeito ou a causa?}

\begin{resumo}
Auditoria pública é disciplina de inferência controlada. O auditor parte de critérios, seleciona procedimentos capazes de produzir evidência, avalia suficiência e então conclui dentro do escopo. Em Auditoria de TI, essa lógica é aplicada a acessos, mudanças, operações, aplicações, bancos, nuvem, contratos, dados e IA. Ferramentas técnicas ampliam cobertura, mas não substituem julgamento, validação de dados e cadeia probatória.
\end{resumo}
